%%%%%%%%%%%%
%
% $Autor: Wings $
% $Datum: 2019-03-05 08:03:15Z $
% $Pfad: SafetyGuidelines.tex $
% $Version: 4250 $
% !TeX spellcheck = en_GB/de_DE
% !TeX encoding = utf8
% !TeX root = manual 
% !TeX TXS-program:bibliography = txs:///biber
%
%%%%%%%%%%%%

\chapter{Sicherheitsleitfaden und Benutzerhinweise}
\textbf{WICHTIG:} Lesen Sie diese Anleitung vor der ersten Inbetriebnahme vollständig durch. Bewahren Sie dieses Dokument für die zukünftige Nutzung sorgfältig auf.

\section{Gefahren für den Anwender}
\begin{itemize}
	\item 
	Der Einfluss von Alkohol, Medikamenten oder Drogen kann zu diversen Gefahren bei der Verwendung des Produkts führen. Daher darf das Tischförderband nur im nüchternen und gesundheitlich uneingeschränkten Zustand bedient werden.
	\item 
	Falls das Produkt sichtbare Schäden aufweist, so ist eine weitere Verwendung des Produkts ausdrücklich zu unterlassen. Abhängig vom Produkt und von der Beschädigung besteht nicht nur Verletzungsgefahr (z.B. durch scharfe Kanten), sondern auch Lebensgefahr (z.B. durch einen elektrischen Spannungsübertritt).
	
	\end{itemize}

\section{Stromversorgung und Netzteil}
\begin{itemize}
	\item \textbf {Netzspannung:}	 
	Verwenden Sie das Netzteil nur mit der korrekten Netzspannung (100 – 240 V Wechselstrom). Schließen Sie es an eine gut erreichbare Steckdose in unmittelbarer Nähe des Geräts an.	
	\item \textbf {Einsatzort:}
	Das Netzteil ist nur für den Betrieb in trockenen Innenräumen geeignet. Schützen Sie es vor Feuchtigkeit und berühren Sie es niemals mit nassen oder fettigen Händen.
	\item \textbf {Aufsicht:}
	Ein Gebrauch durch minderjährige Benutzer in Eigenverantwortung ist nicht gestattet. Unter Aufsicht volljähriger Personen, welche den entsprechenden Sicherheitsanweisungen des Tischförderbands kundig sind, ist eine Benutzung an dieser Stelle zulässig.  \\
	\item \textbf {Beschädigungen:}
	Trennen Sie das Netzteil sofort vom Strom, wenn Sie ungewöhnliche Geräusche, Gerüche oder Rauchentwicklung bemerken. Zerlegen Sie das Netzteil nicht. Bei Beschädigungen (z. B. am Stecker oder Gehäuse) muss das Netzteil entsorgt und durch ein Original-Ersatzteil ersetzt werden.
	\item \textbf {Gewitter:}
	Berühren Sie das Gerät oder das Netzteil während eines Gewitters nicht, wenn es am Stromnetz angeschlossen ist.
\end{itemize}


\section{Mechanische Sicherheit und Aufbau}
\begin{itemize}
	\item \textbf {klappbare Stützbeine:}
	Das Gerät verfügt über klappbare Stützbeine. Prüfen Sie vor jedem Betrieb, ob die klappbaren Stützbeine vollständig ausgeklappt und fest arretiert sind. Ein instabiler Stand kann zum Umkippen der Last oder zum Einklemmen von Gliedmaßen führen.
	\item \textbf {Oberfläche:}
	Stellen Sie das Tischförderband nur auf eine ebene, rutschfeste Fläche.
	\item \textbf {Kleidung und Haare:}
	Tragen Sie keine weite Kleidung, Schals oder Schmuck, die vom Förderband oder den Laufrollen erfasst werden könnten. Langes Haar muss zusammengebunden werden.
\end{itemize}


\section{WARNUNG: Fehlende Schutzeinrichtungen}
\begin{itemize}
	\item \textbf {ACHTUNG:}
	Das Gerät befindet sich in einem unvollständigen Sicherheitszustand, da die geplanten Lichtschranken am Anfang und Ende des Bandes noch nicht installiert sind.
	\item \textbf {Kein Automatik-Stopp:}
	Das Band stoppt aktuell nicht selbstständig. Objekte am Bandende können herunterfallen oder sich verklemmen.
	\item \textbf {Manuelle Überwachung:}
	Das Gerät darf in diesem Zustand nur unter ständiger Aufsicht betrieben werden. Der Bediener muss jederzeit bereit sein, die \glqq Pause"- oder \glqq Stop"-Taste manuell zu drücken.
	\item \textbf {Einzugsgefahr:}
	Achten Sie besonders auf die Umlenkrollen an beiden Enden. Greifen Sie niemals in das laufende Band
\end{itemize}

\section{Wartung und Pflege}
\begin{itemize}
	\item \textbf {Reinigung:}
	Reinigen Sie das Produkt nur im stromlosen Zustand (Netzstecker ziehen). Wischen Sie es mit einem trockenen, weichen Tuch ab. Verwenden Sie keinesfalls Verdünner, Benzin oder aggressive Lösungsmittel, sowie Wasser.
	\item \textbf {Lagerung:}
	Setzen Sie das Tischförderband keinen hohen Temperaturen (durch Feuer oder andere Hitzequellen) oder direkter Sonneneinstrahlung aus. Lagern Sie das Tischförderband in trockener Umgebung bei Raumtemperatur. 
	\item \textbf {Kleinteile:}
	Bewahren Sie Verpackungsmaterialien und lose Gegenstände wie das USB-kabel außerhalb der Reichweite von Kleinkindern auf (Erstickungs- und Strangulationsgefahr).
\end{itemize}

\section{Entsorgung}
\begin{itemize}
	\item 
	Entsorgen Sie dieses Produkt nicht im Hausmüll. Da es sich bei dem Tischförderband um ein elektronisches Gerät handelt, ist dieses an entsprechender Stelle (z.B. kommunaler Wertstoffhof) zu entsorgen. 
	\item \textbf {Rückgabe:}
	In Deutschland sind Vertreiber von Elektrogeräten verpflichtet, Altgeräte unter bestimmten Bedingungen unentgeltlich zurückzunehmen. Adressen kostenfreier Sammelstellen erhalten Sie bei Ihrer Kommunalverwaltung.
	\item \textbf {Datenschutz:}
	Sollte das Gerät personenbezogene Daten enthalten, sind Sie selbst für deren Löschung verantwortlich.
\end{itemize}